\documentclass[10pt,letterpaper]{article}
\usepackage{cloudsuite}
\renewcommand{\DocShort}{CH. 15}
\begin{document}
\coverpage{Chapter 15: Specialized Cloud Architectures}{Cornell study notes and applied review | ARCHITECTURES}{Concepts \textbullet{} decisions \textbullet{} verification artifacts}
\sourceStatement
\tableofcontents
\newpage
\section{Learning contract}
\begin{takeaway}
Choose specialist patterns only when their operating constraints justify them.
\end{takeaway}
\textbf{Chapter scope:} I/O, tiering, networks, edge/fog, abstraction and federation.\\
\sourceRef{15}{415}
\subsection{Mapped source sections}
\begin{itemize}
\item 15.1 Direct I/O Access Architecture
\item 15.2 Direct LUN Access Architecture
\item 15.3 Dynamic Data Normalization Architecture
\item 15.4 Elastic Network Capacity Architecture
\item 15.5 Cross-Storage Device Vertical Tiering Architecture
\item 15.6 Intra-Storage Device Vertical Data Tiering Architecture
\item 15.7 Load-Balanced Virtual Switches Architecture
\item 15.8 Multipath Resource Access Architecture
\item 15.9 Persistent Virtual Network Configuration Architecture
\item 15.10 Redundant Physical Connection for Virtual Servers Architecture
\item 15.11 Storage Maintenance Window Architecture
\item 15.12 Edge Computing Architecture
\item 15.13 Fog Computing Architecture
\item 15.14 Virtual Data Abstraction Architecture
\item 15.15 Metacloud Architecture
\item 15.16 Federated Cloud Application Architecture
\end{itemize}
\section{Cornell cue and notes matrix}
\small\setlength{\LTleft}{0pt}\setlength{\LTright}{0pt}
\begin{longtable}{>{\RaggedRight\arraybackslash}p{0.29\textwidth}>{\RaggedRight\arraybackslash}p{0.63\textwidth}}
\toprule\textbf{Cue / essential question}&\textbf{Explanatory notes and reasoning}\\\midrule\endfirsthead\toprule\textbf{Cue}&\textbf{Notes}\\\midrule\endhead
\textbf{Why use direct I/O and LUN access?} & Direct access may reduce abstraction overhead but can limit portability and complicate ownership, contention or isolation.\\[3pt]\textcolor{petrol}{\emph{Practice:} Compare baseline performance and operational complexity.}\\[8pt]\midrule
\textbf{What is storage tiering optimizing?} & Cross-device and within-device tiering place data on performance or cost classes according to observed access or policy.\\[3pt]\textcolor{petrol}{\emph{Practice:} Define a hot-data promotion rule and prove it works.}\\[8pt]\midrule
\textbf{What is the objective of multipath or redundant connections?} & Alternative network/storage paths reduce the risk that one broken link interrupts access, provided paths avoid a shared failure.\\[3pt]\textcolor{petrol}{\emph{Practice:} Failure-test independent links.}\\[8pt]\midrule
\textbf{How do virtual switches and persistent configurations help?} & Virtual networking can balance work and preserve logical identity as virtual resources are moved or recreated.\\[3pt]\textcolor{petrol}{\emph{Practice:} Migrate a workload and inspect route/policy continuity.}\\[8pt]\midrule
\textbf{Where do edge and fog computing fit?} & Edge processing occurs close to data production; fog resources add intermediary processing tiers between endpoints and centralized cloud components.\\[3pt]\textcolor{petrol}{\emph{Practice:} Map a sensor-to-cloud data flow and a disconnected mode.}\\[8pt]\midrule
\textbf{What do data abstraction, metacloud and federation change?} & These models span storage representations, cloud control coordination or collaborating cloud applications. They require explicit semantics and cross-domain governance.\\[3pt]\textcolor{petrol}{\emph{Practice:} Draw three diagrams showing data, management and application integration separately.}\\[8pt]\midrule
\end{longtable}\normalsize
\section{Synthesis and architecture connections}
An architecture name is a repeatable composition of mechanisms with operating assumptions. Draw component relationships, failure domains, trust boundaries and a representative workload before committing to the pattern.
\begin{workbox}[title={Applied deliverable}]
Choose specialist patterns only when their operating constraints justify them.\\[4pt]Create a one-page architecture or operations artifact with: (1) assumptions, (2) a decision and alternatives, (3) control ownership, (4) observable acceptance criteria, and (5) a failed-case test.
\end{workbox}
\section{Active recall}
\begin{enumerate}
\item Why use direct I/O and LUN access?
\item What is storage tiering optimizing?
\item What is the objective of multipath or redundant connections?
\item How do virtual switches and persistent configurations help?
\end{enumerate}
\subsection{Answer-check rubric}
\begin{itemize}
\item A correct response distinguishes the named concepts rather than relying on vendor names.
\item The explanation identifies an operational consequence, boundary or tradeoff.
\item At least one measurement, test or documentary artifact could demonstrate the claim.
\end{itemize}
\section{Blank Cornell capture sheet}
\textbf{Topic / source subsection:}\enspace\rule{0.65\linewidth}{0.4pt}\\[6pt]
\begin{longtable}{p{0.28\textwidth}p{0.64\textwidth}}
\toprule Cues / questions & Notes / evidence \\ \midrule
\rule{0pt}{42pt} & \\ \midrule
\rule{0pt}{42pt} & \\ \midrule
\rule{0pt}{42pt} & \\ \midrule
\rule{0pt}{42pt} & \\ \midrule
\rule{0pt}{42pt} & \\ \midrule
\end{longtable}
\textbf{Summary in three sentences:}\\[16pt]\hrulefill\\[16pt]\hrulefill\\[16pt]\hrulefill

\refs
\end{document}
